\begin{table*}[t]
\centering
\small
\setlength{\tabcolsep}{2.5pt}
\renewcommand{\arraystretch}{1.16}
\begin{tabular*}{\textwidth}{@{\extracolsep{\fill}}p{0.04\textwidth}p{0.25\textwidth}p{0.65\textwidth}@{}}
\toprule
\textbf{RQ} & \textbf{Memory capability} & \textbf{Benchmark and observable prediction} \\
\midrule
1 & Step-by-step source reading & MemoryAgentBench compares direct retrieval, repeated read-only memory calls, and source reading with citation checks on the same 767 questions. Scores should separate most on fact consolidation and multi-mention tasks when the agent can keep reading. \\
2 & Temporal and multi-session access & LongMemEval tests updates and cross-session links. Versioned observations should keep both earlier and later facts available, while preference questions remain a boundary for latent-user-state modeling. \\
3 & Relational consolidation & LoCoMo separates single-hop, multi-hop, temporal, and open-domain questions. Concept families and relation links should matter most when an answer joins episodes or revises an earlier value. \\
4 & Revisable construction & Base/Align+ paired runs measure duplicate families, ingest calls, tokens, and task slices. Progressive alignment predicts lower construction cost and better consolidation without requiring a uniform score increase. \\
5 & Deletion and absence boundary & MEME tests deletion, cascade, and absence. The current append-only store should retain source links but cannot yet guarantee suppression after deletion. \\
\bottomrule
\end{tabular*}
\caption{Research questions and measurements. Each row tests a distinct prediction of the search--evidence separation; published comparison rows keep their test details so the score differences can be checked.}
\label{tab:experiment-map}
\end{table*}


